%=============================================================================!
% 15.10  WHAT THE TRAJECTORY LOOKS LIKE                                       !
%=============================================================================!
\section{What the trajectory looks like}
\label{sec:cv-trajectory}

Each fault of this chapter lies in the numbers drawn from a trajectory,
and each shows on one of four plots of the quantity concerned, which
together we call its dashboard: the series with the start found, the
running mean from that start with its error band, the errors from blocks,
and the autocorrelation function. This section shows the dashboard of a
healthy run, then six runs or analyses that go wrong, each with the check
that catches it.

\subsection*{A healthy run}

\Cref{fig:cv-healthy} is the dashboard of the volume in the
\SI{1}{\nano\second} at constant pressure that ends the protocol of
\cref{sec:cv-protocol}, drawn by \texttt{dashboard} of the chapter's
scripts. (a) The series has settled by the start found at
\SI{10}{\pico\second} and then wanders about its mean with no trend. (b)
The running mean from the start swings while few independent samples are
in it and then holds steady inside its error band, which narrows as one
over the square root of the time since the start. (c) The error from
blocks rises to a plateau, 9.6 with its own error of 1 against
\SI{9.8}{\angstrom\cubed} from $g$. (d) The autocorrelation function falls
to zero by \SI{8}{\pico\second}, where the estimate of $g$ stops. The
report, given no target, passes all four checks: $V = (12310.4
\pm 9.8)$\,\si{\angstrom\cubed} from 316 independent samples, with a slope
0.85 of its error.

\begin{figure}[htb]
\centering
\includegraphics{ch15_convergence/healthy}
\caption{The dashboard of a converged quantity: the volume over
\SI{1}{\nano\second} at \SI{0.1}{\giga\pascal}. (a) The series, the start
found (dotted) and the mean after it (dashed). (b) The running mean from
the start with its error band. (c) The error of the mean from blocks,
against $\sqrt{g s^2/n}$ (dashed). (d) The autocorrelation function, with
the lag at which the estimate of $g$ stops (dotted).}
\label{fig:cv-healthy}
\end{figure}

\subsection*{Six faults}

\Cref{fig:cv-faults} and \cref{tab:cv-faults} set six faults against
healthy runs.

\begin{table}[htb]
\centering
\small
\caption{The faults of \cref{fig:cv-faults}, the check that catches each,
and its value in the healthy and the faulty run or analysis.}
\label{tab:cv-faults}
\begin{tabular}{@{}>{\raggedright\arraybackslash}p{0.22\linewidth}
   >{\raggedright\arraybackslash}p{0.27\linewidth}
   >{\raggedright\arraybackslash}p{0.21\linewidth}
   >{\raggedright\arraybackslash}p{0.21\linewidth}@{}}
\toprule
fault & check & healthy & faulty\\
\midrule
(a) averaged from the start & the start found; the mean after it &
start \SI{2.8}{\pico\second}, $\bar U/N = (-50.060 \pm
0.085)$\,\si{\milli\electronvolt} &
$\bar U/N = \SI{-50.173}{\milli\electronvolt}$ from $t = 0$\\
(b) the naive error bar & pieces whose bars hold the mean of
\SI{2}{\nano\second} & 11 of 20 with $\sqrt g$ & 0 of 20 with $s/\sqrt
n$\\
(c) too short for a plateau & independent samples; the blocks &
1369 in \SI{2}{\nano\second} & 19.7 in \SI{20}{\pico\second}, no
plateau\\
(d) a drifting energy & the drift test & $-0.9$ errors at
\SI{10}{\femto\second} & $+9.3$ errors at \SI{35}{\femto\second}\\
(e) the wrong ensemble & slope of the ratio of densities, over the
required & $1.05 \pm 0.06$ (CSVR) & $5.50 \pm 0.55$ (Berendsen)\\
(f) a box too small & $\ens{\Delta r^2}/6t$ at \SI{20}{\pico\second}, in
\SI{e-4}{\angstrom\squared\per\femto\second} &
4.177 for 2048 atoms & 3.863 for 256\\
\bottomrule
\end{tabular}
\end{table}

\emph{(a) Averaged from the start.} A run started on the lattice and
averaged from its first step includes the melting: the mean of $U/N$ over
the \SI{100}{\pico\second} is \SI{-50.173}{\milli\electronvolt}, while from
the start found at \SI{2.8}{\pico\second} it is $(-50.060 \pm
0.085)$\,\si{\milli\electronvolt}, against \SI{-50.092}{\milli\electronvolt}
over \SI{2}{\nano\second}. The running mean from $t = 0$ in
\cref{fig:cv-faults}(a) creeps up for the whole run, the offset fading
only as one over the length of the run. In this run the shift is about one
error bar; over the eight runs of \cref{sec:cv-equilibration} the discard
moves the mean by \SI{0.106}{\milli\electronvolt}, 1.6 of their
error bars. Even the \SI{2}{\nano\second} run, started from an
equilibrated liquid, has its start found at \SI{30}{\pico\second}: the
search costs little where there is nothing to discard.

\emph{(b) The naive error bar.} $s/\sqrt n$ treats the 10\,000 samples of
a piece of \SI{100}{\pico\second} as independent. Its bars are too short
by the factor $\sqrt g$, about twelve, and none of the twenty holds the
mean of the whole run, where the bars with $\sqrt g$ hold it in eleven,
against the fourteen that 68\% would give, within the spread of two that
twenty pieces allow (\cref{fig:cv-faults}(b)). The check is to report $g$ with every error
bar.

\emph{(c) Too short for a plateau.} The first \SI{20}{\pico\second} of the
run hold 19.7 independent samples of $U$. Its block errors rise without
levelling (\cref{fig:cv-faults}(c)), and its start falls in the second
half of the run, at \SI{12}{\pico\second}: the report fails it on the start
and on 100 independent samples. Its block check passes, because the
\SI{8}{\pico\second} after that start give $g = 13$, a tenth of the value
of the whole run, and blocks of $5g$ that short agree with an error as
short; the plot of the blocks shows what the check cannot.

\emph{(d) A drifting energy.} With steps of \SI{35}{\femto\second} the
total energy of a run at fixed energy climbs by
\SI{0.6}{\milli\electronvolt} per atom in a nanosecond
(\cref{fig:cv-faults}(d)), 9.3 errors of the drift test, and moves the
temperature by \SI{3.2}{\kelvin} per nanosecond; with steps of
\SI{10}{\femto\second} its slope is $-0.9$ errors, no drift.

\emph{(e) The wrong ensemble.} Berendsen coupling with $\tau_{\mathrm T} =
\SI{0.1}{\pico\second}$ gives the right mean energies and the wrong
spread, and the logarithm of the ratio of its energy densities at 130 and
\SI{140}{\kelvin} rises 5.5 times too steeply (\cref{fig:cv-faults}(e)).
None of the convergence checks finds this, since each run is converged in
its own narrow ensemble; the test of \cref{sec:cv-ensemble} does, as does
the spread of $K$ set against the canonical one (\cref{sec:th-compare}).

\emph{(f) A box too small.} The squared displacement over $6t$, whose
value at long times is $D$, reaches \SI{3.863e-4}{\angstrom\squared\per
\femto\second} at \SI{20}{\pico\second} for 256 atoms and $4.177\times
10^{-4}$ for 2048, so the smaller box gives 0.925 times the larger's
(\cref{fig:cv-faults}(f)); these are single runs read at one time, near
the 3.914 and 4.214 of \cref{sec:cv-size}. Each run fixes its own $D$
to about 1\%; the smaller box fixes it at a different number, for the
reason \cref{sec:cv-size} gives.

\begin{figure}[htb]
\centering
\includegraphics{ch15_convergence/faults}
\caption{Faults of convergence (teal) against healthy runs (grey dashed).
(a) The running mean of $U/N$ of a melting run from $t = 0$, against that
from the start found, and the mean of \SI{2}{\nano\second} (dotted). (b)
Twenty pieces of \SI{100}{\pico\second}: the mean of $U$ less that of the
whole run, with $s/\sqrt n$ (teal, too short to see) and $\sqrt{g s^2/n}$
(grey). (c) The error of $U$ from blocks over the first
\SI{20}{\pico\second} and over \SI{2}{\nano\second}, each over its
$\sqrt{g s^2/n}$. (d) The total energy per atom at fixed energy at 35 and
\SI{10}{\femto\second}. (e) The logarithm of the ratio of the energy
densities at 140 and \SI{130}{\kelvin} under Berendsen coupling (here in
teal, as the fault) and CSVR, each about its mean. (f) The squared displacement over $6t$ for 256 and
2048 atoms.}
\label{fig:cv-faults}
\end{figure}

\begin{practice}
A run is read through the dashboard of each quantity to be reported: the
series with its start, the running mean with its error band, the block
errors with their plateau, and the autocorrelation function. Every average
is reported with its error bar, its number of independent samples and the
start discarded, and passes \texttt{convergence\_report}; a run that fails
is extended by the length \cref{eq:cv-length} gives. A thermostat is
checked once by the test of \cref{sec:cv-ensemble}, a time step by the
drift test at fixed energy, and a transport coefficient by a second box
size.
\end{practice}

\begin{keyidea}
Each fault of convergence shows on a plot of the dashboard or on one of the
chapter's tests: the start for an average taken from the first step, $g$ for
a naive error bar, the block plateau for a run too short, the drift test
for a creeping energy, the ratio of densities for a thermostat that
narrows the spread, and a second box for a size-dependent result.
\end{keyidea}

\notebook{15}{break each run in turn and watch which check catches it}

\begin{selfcheck}
A paper reports $D$ from one run of 256 argon atoms as $(3.87 \pm 0.04)
\times10^{-4}$\,\si{\angstrom\squared\per\femto\second}. Is the error bar
the error of $D$?
\end{selfcheck}
